Après l'analyse syntaxique, on obtient une représentation structurée du programme, sous la fore d'un arbre syntaxique abstrait (AST). Cependant, cette représentation ne garantit pas que le programme est sémantiquement correct. 


\begin{exemple}{}{}
Le programme peut mutliplier un entier par un flottant, ce qui est une opération soit à invalider, soit à convertir implicitement. L'analyse sémantique va donc vérifier la validité des opérations et effectuer les conversions nécessaires.
\end{exemple}

Il y a donc au cours de cette phase de l'analyse sémantique, deux tâches principales~:
\begin{itemize}
  \item analyse de type~: vérifier que les opérations sont valides pour les types d'opérandes donnés, et effectuer les conversions nécessaires si possible.
  \item analyse de noms~: vérifier que les identificateurs utilisés dans le programme sont déclarés et accessibles dans le contexte actuel.
\end{itemize}
